% Lösungen Einheit 3
\einheitenkopf*{Einheit 3 von 3 · Kreisteile}
\erg{20}{a) $\frac12$ \quad b) $\frac14$ \quad c) $\frac13$ \quad d) $\frac34$ \quad e) $\frac16$}
\erg{21}{a) $\frac14$ \quad b) $\frac12$ \quad c) $\frac16$ \quad d) $\frac18$ \quad e) 10\,\% \quad f) $50 : 360 \approx 13{,}9$\,\%}
\erg{22}{a) 45° \quad b) 135° \quad c) $0{,}3 \cdot 360° = 108°$}
\erg{23}{a) 11,0\,cm \quad b) 6,1\,cm \quad c) 21,4\,cm² \quad d) $6{,}11 + 2 \cdot 7 \approx 20{,}1$\,cm \quad e) $b \approx 17{,}3$\,m, $A \approx 64{,}8$\,m²}
\erg{24}{a) $\pi \cdot (64 - 36) \approx 88{,}0$\,cm² \quad b) $\pi \cdot (121 - 1) \approx 377{,}0$\,m² \quad c) $r_a = 11$\,cm, $r_i = 8$\,cm, $A = \pi \cdot (121 - 64) \approx 179{,}1$\,cm² \quad d) nicht mit dem Bus: $360° - 125° = 235°$; $235 : 360 \approx 65{,}3$\,\%}
\erg{25}{a) Ben teilt durch 180° statt durch den Vollwinkel 360°. Richtig: $150 : 360 \approx 0{,}417 = 41{,}7$\,\% \quad b) $162 : 360 = 0{,}45 = 45$\,\%}
\erg{26}{a) Ein Halbkreis ist die Hälfte des Kreises, also gehört die Hälfte des Vollwinkels dazu: $360° : 2 = 180°$; die beiden Radien bilden zusammen einen Durchmesser. \quad b) Ja: Der Anteil $\alpha : 360°$ verdoppelt sich, der Radius bleibt gleich. \quad c) Nein: Es kommt auch auf den Radius an; ein 60°-Ausschnitt aus einem großen Kreis kann größer sein.}
\erg{27}{a) $\frac14 \cdot \pi \cdot 9^2 \approx 63{,}6$\,m² \quad b) $130 : (\pi \cdot 81) \cdot 360° \approx 183{,}9°$, also etwa 184° \quad c) Nein: 184° ist mehr als 180°; bei 180° bewässert er nur $\frac12 \cdot \pi \cdot 81 \approx 127{,}2$\,m², etwas zu wenig.}
